% ============================================================
% Introduction — synthesis of P1-P3 (target: Computers & Security)
% ============================================================
\section{Introduction}
\label{sec:intro}

Phishing detection is often presented as an almost solved problem. Modern email
filters, built on the \emph{content} signal (lexical features, URLs, attachments) and reputation,
achieve very high effectiveness on mass campaigns --- on the order of $\sim$98\%
correct classifications on typical benchmarks \citep{birthriya2025detection,tan2020graph}. This number
is, however, misleading. The remaining fraction of traffic --- \textbf{spear phishing} and \textbf{Business Email
Compromise (BEC)} --- accounts for a disproportionately large share of financial losses
\citep{cidon2019high,ho2019detecting}, while at the same time being almost undetectable by content-based methods.
\emph{It is precisely these ``$2\%$'' that are the real problem.}

\paragraph{Why the residual $2\%$ gets through.} We argue that this is neither accidental nor a
transient gap, but a \emph{structural} property of the targeted attack. Spear phishing succeeds
\textbf{precisely because} it \emph{respects the domain knowledge} of the victim: their roles
and hierarchy, business relationships, competencies, communication routines, and project context. A request
from the ``CEO'' to the CFO for an urgent transfer is \emph{content-wise} indistinguishable from
a legitimate one --- it is \emph{plausible} in light of knowledge about the organization \citep{butavicius2022why,sturman2025security}.
The attacker builds this knowledge from OSINT \citep{dewan2014analyzing,brown2025web}, and language
models drastically reduce its cost \citep{bethany2024evaluating}. Since maliciousness is determined not by
\emph{what} is written, but by \emph{who} sent it and \emph{in what context}, the
detection signal must lie \emph{beyond} the content --- in the structure of relationships and in the dynamics of their violation.

\paragraph{Thesis: spear phishing is a violation of domain knowledge.} We formalize the \emph{domain knowledge graph}
of an organization as a multiplex network of relations over a set of accounts (contacts, collaboration, competencies,
OSINT traces, temporal rhythm). In this framing, the targeted attack reveals itself through two invariants,
both invisible to a content filter. The first is \textbf{cross-layer (in)consistency}: a genuine relation is
consistent across many layers at once, whereas impersonation --- external or a compromised account --- breaks
the pattern, appearing ``known'' in one layer yet inconsistent in another (off-rhythm, or without an OSINT
trace). The second is \textbf{propagation dynamics}: lateral phishing is a cascade process on the graph, in which
a compromised account sends to its contacts off its own rhythm and leaves a wave-like signature on the
\emph{receiving} side that a single message does not reveal. The framing covers \emph{uniformly} both modes of
the targeted attack: \textbf{lateral/BEC} (a compromised internal account, a known contact) and \textbf{external
impersonation} of a role or authority.

\paragraph{Research questions and sub-theses.} We make the thesis testable by splitting it into three research
questions, each paired with a measurable sub-thesis that the results are held to.
\begin{description}[leftmargin=1.4em,noitemsep,topsep=2pt]
  \item[\textbf{RQ1} (static signal).] Does cross-layer inconsistency in the domain knowledge graph separate
  targeted attacks from legitimate new senders? \emph{Sub-thesis~T1}: a multiplex of domain-knowledge layers
  raises sensitivity@FPR over a contacts-only graph and removes new-sender false alarms, and the gain vanishes
  under a layer/rhythm shuffle control.
  \item[\textbf{RQ2} (temporal signal).] Can a compromised account --- consistent across all static layers --- be
  detected from cascade dynamics, and does the signal survive a stealthy attacker? \emph{Sub-thesis~T2}: a
  temporal detector reading the receiving-side wave beats volume detection at low FPR and, unlike it, does not
  break down under a stealthy attacker; the advantage vanishes under a time-shuffle control and transfers to real
  topologies.
  \item[\textbf{RQ3} (measurement and limits).] Which organizational roles and relations are most exposed, and
  what are the operational limits of the signal? \emph{Sub-thesis~T3}: susceptibility is governed by relation,
  not rank, and at a realistic base rate the detectors are a triage instrument rather than an autonomous defense.
\end{description}
\noindent The goal is not to deliver a deployable filter or a new learning architecture, but to establish that
the two invariants carry a \emph{causal, transferable} signal and to bound where it stops being operationally
useful; the learning models are off-the-shelf instruments throughout.

\paragraph{Measurement rigor: what it means for a method to ``work''.} This work places particular emphasis on
\emph{research methodology}. We show that the results of graph-based detectors are easy to
\emph{inflate}: real-world anchors introduce \emph{leaks} (of rhythm, degree, community alignment
with labels), and synthetic generators favor certain classes of models. Following
the caution of \citet{arp2022dos} and the TESSERACT protocol \citep{pendlebury2019tesseract},
we introduce a \emph{leak-aware} evaluation and a \emph{predictivity probe}
measuring whether the ranking of methods transfers across topologies. By this same
measure we \emph{define the notion of effectiveness} of a method for residual spear phishing.
Throughout, the learning models we employ --- a temporal graph network and gradient boosting ---
serve as \emph{off-the-shelf measurement instruments} that quantify the two invariants; the
contribution of this work is conceptual and methodological --- a way of \emph{framing and measuring}
targeted attacks --- not a new machine-learning architecture.

\paragraph{Contributions.} The contributions map onto the research questions as follows.
\begin{enumerate}[noitemsep,topsep=2pt]
  \item \textbf{The notion and formalization of the domain knowledge graph} and the thesis that spear phishing is
  its violation --- uniformly for lateral/BEC and external impersonation (the Model, Section~\ref{sec:model});
  the foundation for RQ1--RQ3.
  \item \textbf{Two detection invariants} (cross-layer inconsistency, cascade dynamics) beating strong baselines
  under a content-evasion regime, with causal controls (\textbf{RQ1}, \textbf{RQ2}; T1, T2).
  \item \textbf{An organizational susceptibility metric} measuring which roles, relations, and communication
  patterns are most exposed to spear phishing (\textbf{RQ3}; T3, Section~\ref{sec:susceptibility}).
  \item \textbf{A leak-aware evaluation methodology} and an \emph{operational} definition of detection
  effectiveness, with an honest delineation of the limits --- FPR, base rate, adversarial cost (\textbf{RQ3}; T3).
\end{enumerate}

\paragraph{Organization.} Section~\ref{sec:related} discusses spear phishing, human susceptibility, graph-based
approaches, and evaluation methodology. Section~\ref{sec:model} formalizes the domain knowledge graph, the threat
model, the two invariants, and the definition of effectiveness; Section~\ref{sec:method} gives the verification
procedure and the leak-aware protocol. Section~\ref{sec:results} instantiates the model in a case study and
reports the detection results, the susceptibility measurement, and the rigor analysis, followed by a discussion
of the limits and practical implications.
